\section{Decode a signed halfword at any byte address}\label{pcm-encoding}
The low byte and high byte first form an unsigned pattern
\[
  u=\text{low}+256\,\text{high}.
\]
Interpret it as $u$ when $u<32768$, otherwise as $u-65536$. Thus stored bytes
\texttt{ff 7f} represent 32767, \texttt{00 80} represent minus 32768, and
\texttt{ff ff} represent minus one. Byte order chooses the pattern; signed
interpretation chooses its integer value.

The scalar path loads each byte unsigned, shifts the high byte by eight and ORs
them together. \texttt{seh} then sign extends the low sixteen bits to the native
register width. Without that sign extension, a negative sample would instead
participate as a large positive number. A halfword load would also interpret a
signed value, but an unconditional halfword access does not satisfy this exercise's
arbitrary byte-alignment requirement.

After both inputs are sign extended, their exact sum lies from minus 65536 through
65534. It fits in a signed word. Word addition can therefore retain the exact
intermediate before the output is narrowed. The low sixteen bits alone do not
contain enough information to decide whether clipping occurred.
